% Generated draft: requires mathematical and source review.
\documentclass[11pt]{article}
\usepackage[margin=25mm]{geometry}
\usepackage{fontspec}
\setmainfont{Latin Modern Roman}
\usepackage{amsmath,amssymb,mathtools,unicode-math}
\setmathfont{Latin Modern Math}
\usepackage{enumitem,needspace,xurl,hyperref}
\setlength{\parindent}{0pt}
\setlength{\parskip}{5pt}
\setlength{\emergencystretch}{3em}
\newcommand{\N}{\mathbb{N}}
\newcommand{\Z}{\mathbb{Z}}
\newcommand{\Q}{\mathbb{Q}}
\newcommand{\R}{\mathbb{R}}
\newcommand{\C}{\mathbb{C}}
\newcommand{\F}{\mathbb{F}}
\newcommand{\A}{\mathbb{A}}
\newcommand{\PP}{\mathbb{P}}
\newcommand{\E}{\mathbb{E}}
\newcommand{\eps}{\varepsilon}
\newcommand{\dd}{\,\mathrm{d}}
\begin{document}
\section*{UnsolvedMath ID: 30004411}
\textbf{Source problem code:} catalogue entry
\subsection{Definitions and mathematical statement}
For a smooth knot \(K\subset S^3=\partial B^4\), define its zero-framed trace by
\[
X_0(K)=B^4\cup_{K,0}(D^2\times D^2),
\]
where a smooth \(2\)-handle is attached along \(K\) with framing zero relative to the Seifert longitude. The supplied record writes \(X(K)\) without defining its framing. The formulations below use the conventional interpretation \(X(K)=X_0(K)\), subject to confirmation from the report. A knot is smoothly slice if it bounds a smoothly and properly embedded disk in the standard smooth \(B^4\). The word ``nonslice'' below is interpreted as not smoothly slice; the supplied statement does not explicitly distinguish smooth from topological sliceness.

\begin{enumerate}
\item Let \(K_C\) be the Conway knot. The first question concerns the expression
\[
C_f=S^4\setminus\nu\bigl(f(X(K_C))\bigr)
\]
and asks whether \(C_f\) admits a smooth structure. Here \(f\) denotes the embedding referred to in the source, and \(\nu\) denotes its indicated neighborhood. The supplied record does not specify \(f\), its embedding category, its construction, or the neighborhood convention. Consequently it does not determine a particular complement. In particular, it leaves unresolved whether the intended complement is the compact exterior \(S^4\setminus\operatorname{int}N\), for a closed neighborhood \(N\), or a complement with its boundary removed. The accompanying supplied assessment refers to a proposed topological four-sphere construction. This question must therefore not be read as assuming that \(f\) is already a smooth embedding into the standard smooth \(S^4\). The record also does not state whether a desired smoothing must respect any specified boundary smoothing or other prescribed data.

\item Does there exist a smooth integer-homology \(4\)-sphere \(M\), together with a smooth knot \(K\subset S^3\) that is not smoothly slice, and a smooth embedding
\[
\iota:X(K)\hookrightarrow M?
\]
Here an integer-homology \(4\)-sphere means a closed, connected, oriented smooth \(4\)-manifold satisfying
\[
H_i(M;\Z)\cong H_i(S^4;\Z)\qquad\text{for every }i\geq 0.
\]
No simple-connectivity assumption on \(M\) is included. This follows the supplied clean formulation's use of ``integer-homology \(4\)-sphere''; the original wording only says a smooth \(4\)-manifold with the integer-homology type of \(S^4\), so the intended closedness should be confirmed. The wording ``of a nonslice knot'' suggests an existential choice of \(K\), although the preceding question fixes the Conway knot; that scope also requires confirmation.

\item For a specified smooth knot \(K\subset S^3\), consider the following property: there exist closed, oriented, simply connected smooth \(4\)-manifolds \(X\) and \(W\) such that
\[
X\simeq W,\qquad
\exists\,\iota:X(K)\hookrightarrow X\ \text{smoothly},\qquad
\nexists\,j:X(K)\hookrightarrow W\ \text{smoothly}.
\]
Here \(X\simeq W\) means that the underlying spaces are homotopy equivalent. The source asks whether such manifolds exist, but does not specify whether \(K\) remains the Conway knot, is the nonslice knot from the second question, or is separately existentially quantified. Thus the displayed property is precise for fixed \(K\), while the intended quantifier on \(K\) remains unresolved. No orientation-preserving requirement on the homotopy equivalence or embeddings is stated in the supplied record.
\end{enumerate}

These are three related questions, not an assertion that one common choice of knot or ambient manifold satisfies all three. Their supplied open status and verification labels are curation claims, not an independent mathematical review of this draft.
\subsection{Short English statement}
The record asks three questions about knot traces: whether a specified Conway-knot trace complement can be smoothed, whether a smooth integer-homology \(4\)-sphere can contain the trace of a nonslice knot, and whether homotopy-equivalent simply connected closed smooth \(4\)-manifolds can differ in admitting a smooth embedding of a knot trace. The trace framing, the first embedding and complement, and the scope of the knot variable require clarification.
\subsection{Sources}
\begin{itemize}
\item Local UnsolvedMath registry: permanent ID \(30004411\), source problem code \texttt{catalogue entry}, title ``Smooth Embeddings of Knot Traces in Homology Four-Spheres,'' in the supplied Oberwolfach Reports Open Problems collection. The supplied proposed year is \(2020\); no individual proposer is identified.
\item Supplied report citation: \emph{Low-dimensional Topology (2021)}, Oberwolfach Reports, pp.\ 465--516, DOI \texttt{10.4171/owr/2020/8}, \url{https://doi.org/10.4171/owr/2020/8}. The same link is supplied in the literature sources under ``Oberwolfach Report 8/2020.'' The citation's \(2021\) label and the report's \(2020\) identifier are retained as supplied; their relationship has not been independently checked. The background describes these reports as documenting mathematical workshops and open-problem discussions.
\item Supplied literature source: Piccirillo, \emph{The Conway knot is not slice}, \url{https://arxiv.org/abs/1808.02923}. The supplied assessment attributes to Piccirillo the proof that the Conway knot is not smoothly slice, and cautions that this does not establish whether the associated complement in the proposed topological four-sphere construction admits or lacks a smooth structure.
\item The exact supplied literature-check date is \texttt{2026-08-22}. The supplied status is ``open.'' This is supplied dated curation, described in the record as literature triage derived from the OWR curation, and has not been independently verified here. Time-sensitive status and cited claims require independent verification.
\item The supplied statement-status label is \texttt{corrected\_verified}. Its verification note reports that the official report was checked and all three listed open problems about knot traces and smooth four-manifolds were restored. This is a supplied provenance claim; no source visit or independent verification is asserted here.
\end{itemize}
\end{document}
